\documentclass[11pt]{article}
\usepackage{amsmath}
\usepackage{booktabs}
\usepackage[margin=1in]{geometry}
\usepackage{natbib}

\begin{document}

\subsection*{Reinforcement-learning trust-game baseline}

\paragraph{Model.}
As a baseline against the hierarchical AL-RNN, we fit a sequential
reinforcement-learning (RL) model of trust-game investment combining online
learning of trustee/cue repayment ratios with the inequity-averse utility of
\citet{FehrSchmidt1999}, following the formulation in Section~1 of the
supplementary model description. On each trial the predicted repayment
ratio for the presented identity/expression cue is a linear read-out of a
weight vector, $r_t = w_t \cdot \phi_t$, with $\phi_t$ a 9-dimensional (4
identity + 5 expression one-hot) feature vector and $w_0 = \mathbf{0}$ (no
cue priors). Own and trustee gains for investment $a \in \{1,\dots,5\}$
(in units of 10 currency points) are
\begin{equation}
  G_i(a) = (5-a) + 3ar_t, \qquad G_j(a) = 3a - 3ar_t,
\end{equation}
combined into the piecewise-linear Fehr--Schmidt utility
\begin{equation}
  U_t(a) = G_i(a) - \alpha^{FS}\max\!\big(G_j(a)-G_i(a),0\big)
                    - \beta^{FS}\max\!\big(G_i(a)-G_j(a),0\big),
  \qquad \alpha^{FS} \ge \beta^{FS} \ge 0,
\end{equation}
with one $(\alpha^{FS},\beta^{FS})$ pair per trustee. Investments are chosen
via softmax over $U_t$ (inverse temperature $\beta$), and $w_t$ is updated
after the true repayment ratio $\tilde r_t$ is observed by a
prediction-error delta rule, $w_{t+1} = w_t + \alpha(\tilde r_t - r_t)\phi_t$.
Parameters were fit per subject by maximum likelihood on the first 60 of 80
trials -- the identical temporal train/test split used for the AL-RNN --
and scored by mode-prediction MAE and predictive NLL on the held-out final
20 trials.

\paragraph{Extensions tested.}
To determine whether any shortfall relative to the AL-RNN reflects the RL
model's sequential value learning or its decision rule, we tested three
extensions of the base model before ablating the decision rule itself: (1) a
sticky-choice term $\kappa$ added to $U_t(a_{t-1})$, capturing
choice-history bias independent of value; (2) hierarchical fitting via
iterated empirical-Bayes shrinkage, alternating per-subject MAP estimation
under a population Gaussian prior with re-estimation of that prior's
mean/variance across subjects, analogous in spirit to the AL-RNN's
cross-subject pooling; and (3) replacement of the fixed learning rate
$\alpha$ with a per-feature visit-count step size ($1/n$), giving each cue
weight its fastest possible \emph{causal} (no-lookahead) convergence. We
then isolated the decision rule by holding the identical, causally-learned
belief $r_t$ fixed and substituting three alternative decoders in place of
the Fehr--Schmidt utility + softmax: a free softmax regression on $r_t$
(\textsc{softmax-r}), an ordinal cumulative-link model on $r_t$
(\textsc{ordinal-r}), and a free softmax regression directly on the raw cue
identity, bypassing $r_t$ entirely (\textsc{softmax-cue}) -- the latter
included as an upper-bound reference, structurally equivalent to a direct
regression of investment on cue identity.

\paragraph{Justification.}
This staged ablation avoids conflating two distinct failure modes -- slow or
noisy sequential value learning versus a mismatched decision rule -- before
drawing conclusions about the RL framework as a whole. Extensions (1)--(3)
target learning dynamics and left performance essentially unchanged
(Table~\ref{tab:rl_results}), ruling out convergence speed, choice
perseveration, and cross-subject overfitting as explanations. Substituting
only the decoder, while holding $r_t$ fixed, recovered most of the
shortfall, and bypassing $r_t$ entirely (\textsc{softmax-cue}) matched an
independently-fit linear regression of investment on cue identity as well
as the AL-RNN's held-out MAE. We take this as evidence that the specific
inequity-averse decision rule of \citet{FehrSchmidt1999} -- rather than
sequential learning per se -- is responsible for the RL model's shortfall
here. We further note that this trust game draws repayment ratios and cue
identities independently on each trial and resets the investor's endowment
every round, affording little genuine cross-trial dependency beyond a
per-cue additive expectation; this is a likely explanation for why even a
flexible dynamical model does not exceed a static regression's accuracy in
this particular paradigm.

\begin{table}[h]
\centering
\begin{tabular}{lcc}
\toprule
Model & Test MAE & Test NLL \\
\midrule
1. RL baseline (fixed-$\alpha$, FS utility + softmax) & 1.23 & 1.41 \\
2. \quad + sticky-choice $\kappa$                     & 1.26 & 1.41 \\
3. \quad + hierarchical shrinkage                     & 1.23 & 1.41 \\
4. \quad + count-based learning rate                  & 1.23 & 1.46 \\
5. Softmax decoder on $r_t$ (\textsc{softmax-r})       & 0.96 & 1.43 \\
6. Ordinal decoder on $r_t$ (\textsc{ordinal-r})       & 1.01 & 1.40 \\
7. Softmax decoder on raw cue (\textsc{softmax-cue})   & 0.72 & 1.27 \\
\midrule
Linear regression, investment $\sim$ cue (reference)  & 0.8--0.9 & -- \\
AL-RNN (hierarchical)                                  & 0.71 & -- \\
\bottomrule
\end{tabular}
\caption{RL trust-game baseline: mean test-set mode-prediction MAE and
predictive NLL across 32 subjects, 20 held-out trials/subject. Rows 1--4
share the Fehr--Schmidt utility + softmax decision rule and vary only the
value-learning stage; rows 5--7 hold the learned belief $r_t$ fixed (row 7
bypasses it) and vary only the decoder. The linear-regression and AL-RNN
rows are reported from the corresponding analyses elsewhere in the paper
and are included here for reference.}
\label{tab:rl_results}
\end{table}

\subsection*{Bayesian trustee-tracking baseline}
\label{sec:bayes}

\paragraph{Model.}
As a second, structurally distinct baseline, we implemented a discrete
Bayesian trustee-tracking model following \citet{Ray2008}. The investor
maintains, for each of the four trustees $j$, a belief $b_t^{(j)}$ over $K$
discrete ``generosity types'' with fixed return rates $\rho_1 < \cdots <
\rho_K$, initialised uniform. The predicted repayment ratio combines the
belief mean with an expression-cue offset $c(e_t)$:
\begin{equation}
  \widehat{RR}_t = \sum_k b_t^{(j)}(k)\,\rho_k \;+\; c(e_t),
\end{equation}
with $c(e_t)=0$ for the two neutral trustees. Investment values follow
risk-neutral expected-value maximisation, $Q_t(a) = a(3\widehat{RR}_t -
1)$, with choice via softmax, $P(a)\propto\exp(\beta Q_t(a))$. After
observing the true repayment ratio $RR_t$, the shown trustee's belief is
updated by a Bayesian multiply-and-renormalise (additive in log-space, no
learning rate):
\begin{equation}
  b_{t+1}^{(j)}(k) \;\propto\; b_t^{(j)}(k)\,
  \mathcal{N}\!\big(RR_t;\; \rho_k + c(e_t),\ \sigma^2\big).
\end{equation}
The type grid $\rho$ and observation noise $\sigma$ were fixed from the
train-pooled empirical distribution of $RR$ (grid spanning its observed
range; $\sigma$ the residual std after removing trustee-identity and
expression structure), rather than fit through choice likelihood, and
parameters were fit and scored on the identical train/test split and
metrics as the RL baseline.

\paragraph{Structure and complexity variants.}
We probed the same encoder/decoder decomposition as for the RL baseline,
now within the Bayesian model. Two structures from \citet{Ray2008} were
implemented: (B) the full dynamic belief-tracker above, and (C) a static
variant with no trial-by-trial updating, $\widehat{RR} = g_j + c(e_t)$,
with one fixed generosity estimate $g_j$ per trustee. (A third structure,
a low-dimensional shared projection of subject parameters mirroring the
digital-twin architecture, was not pursued: with only $\sim\!5$
interpretable parameters it offers little to compress, and was intended by
\citet{Ray2008}'s formulation mainly for architectural parity rather than
as an independent hypothesis about this task.) Orthogonally, we varied the
expressivity of two stages independently: the \emph{cue-encoding} stage,
where the expression offset $c(e_t)$ was either a single signed slope
(\textsc{linear}: 1 parameter) or a free offset per expression level
(\textsc{free}: 4 parameters, matching the RL model's per-level expression
weights); and the \emph{decoding} stage, where choices were generated
either by the native risk-neutral $Q$-softmax above, or (B$'$) by a free
softmax regression on the identical, causally-tracked belief
$\widehat{RR}_t$, in place of $Q_t(a)$ and $\beta$.

\paragraph{Note on the native decision rule.}
Because $Q_t(a)$ is linear in $a$, its softmax mode is a corner action
($a{=}1$ or $a{=}5$) for any $\beta>0$ and any belief -- risk-neutral
expected-value maximisation cannot produce an interior investment as its
most likely choice, independent of how accurately $\widehat{RR}_t$ is
estimated. This is a property of the decision rule rather than an
estimation failure, and motivates (B$'$) as a fairer test of the
belief-tracking mechanism in isolation.

\begin{table}[h]
\centering
\begin{tabular}{llcc}
\toprule
Model & Cue encoding & Test MAE & Test NLL \\
\midrule
(B) Dynamic + native $Q$-softmax  & linear & 1.15 & 1.35 \\
(B) Dynamic + native $Q$-softmax  & free   & 1.17 & 1.34 \\
(B$'$) Dynamic + flexible decoder & linear & 1.17 & 1.50 \\
(B$'$) Dynamic + flexible decoder & free   & \textbf{0.88} & 1.41 \\
(C) Static, native $Q$-softmax    & linear & 1.08 & 1.29 \\
(C) Static, native $Q$-softmax    & free   & 1.08 & 1.30 \\
\bottomrule
\end{tabular}
\caption{Bayesian trustee-tracking baseline: mean test-set
mode-prediction MAE and predictive NLL across 32 subjects, 20 held-out
trials/subject. ``Cue encoding'' is the expressivity of the
expression-offset term $c(e_t)$ (see above); all other settings (type
grid, $\sigma$, train/test split, evaluation) match the RL baseline
(Table~\ref{tab:rl_results}).}
\label{tab:bayes_results}
\end{table}

\subsection*{Where model complexity pays off}

Across both baselines, performance can be decomposed into an
\emph{encoder} stage that tracks or estimates a belief about the current
trustee/cue (the RL delta-rule weights $r_t$, or the Bayesian tracker's
$\widehat{RR}_t$) and a \emph{decoder} stage that maps that belief onto a
choice (the Fehr--Schmidt utility or risk-neutral $Q$-value, followed by
softmax). In both baselines, increasing the expressivity of the decoder
alone -- replacing a theory-constrained utility/decision rule with a free
regression on the identical, unchanged belief -- produced the largest and
most consistent gains (RL: 1.23 $\to$ 0.96 test MAE, Table~\ref{tab:rl_results};
Bayesian, free cue encoding: 1.17 $\to$ 0.88, Table~\ref{tab:bayes_results}).
Increasing the expressivity of the encoder side alone, without a
correspondingly flexible decoder, contributed comparatively little (the
native-decision-rule rows in Table~\ref{tab:bayes_results} are essentially
unchanged between linear and free cue encoding).

We read this as evidence that, for this task, the two hand-specified
state-tracking mechanisms tested here -- a delta-rule value update paired
with an inequity-averse utility, and a discretised-type Bayesian tracker
paired with risk-neutral expected-value maximisation -- were not, on their
own, sufficiently expressive to capture the structure relating cues to
investment choices, and that substantially more of that structure could be
recovered by relaxing the decoding stage instead. This is a property of
the specific parametric assumptions built into each baseline (a fixed
additive value read-out; a fixed inequity-averse or risk-neutral choice
rule), rather than a general claim that trial-by-trial belief updating is
unnecessary. It also motivates the hierarchical AL-RNN architecture used
elsewhere in this paper: rather than requiring the modeller to decide in
advance how much representational capacity belongs in state-tracking
versus decoding, its encoder and decoder are jointly optimised and can
allocate flexibility to whichever stage the data reward, without that
allocation being specified by hand.

% Bibliography stub -- merge into the paper's existing .bib as needed.
\begin{thebibliography}{9}
\bibitem{FehrSchmidt1999} Fehr, E., \& Schmidt, K. M. (1999). A theory of
  fairness, competition, and cooperation. \textit{The Quarterly Journal of
  Economics}, 114(3), 817--868.
\bibitem{Lee2008} Lee, D. (2008). Game theory and neural basis of social
  decision making. \textit{Nature Neuroscience}, 11(4), 404--409.
\bibitem{RillingSanfey2011} Rilling, J. K., \& Sanfey, A. G. (2011). The
  neuroscience of social decision-making. \textit{Annual Review of
  Psychology}, 62, 23--48.
\bibitem{SuttonBarto2018} Sutton, R. S., \& Barto, A. G. (2018).
  \textit{Reinforcement learning: An introduction}. MIT Press.
\bibitem{Ray2008} Ray, D., King-Casas, B., Montague, P. R., \& Dayan, P.
  (2008). Bayesian model of behaviour in economic games. \textit{Advances
  in Neural Information Processing Systems}, 21 (NIPS).
\end{thebibliography}

\end{document}
